\section{HTTP Flood (DoS)}
\label{sec:threat-dos}

The ops portal is the only service with an unauthenticated public interface and
serves as the operational dashboard for city personnel. A high-rate HTTP flood
against it exhausts its processing capacity, preventing operators from monitoring
sensor readings and responding to environmental alerts. The attack is simulated
using \texttt{hey}, an HTTP load generator, configured with 50 concurrent workers
sustaining requests for 60 seconds, orchestrated by \texttt{scripts/demo\_dos.sh}.

\subsection{Detection}

Detection is implemented as a two-stage Wazuh rule chain on the ops-portal
access logs. Rule 100001 fires on every request to the root or data endpoint
(\texttt{/} or \texttt{/api/data}), producing a low-level marker event. Rule
100002 is a frequency rule that fires when 50 or more such markers arrive from
the same source IP within a 10-second window, confirming an active flood.

\begin{lstlisting}[style=xmlstyle, caption={Wazuh rules for HTTP flood detection.}]
<rule id="100001" level="3">
  <if_sid>86600</if_sid>
  <field name="path">^/$|^/api/data$</field>
  <description>High frequency access to ops-portal
    endpoint</description>
  <group>dos,</group>
</rule>

<rule id="100002" level="12" frequency="50" timeframe="10">
  <if_matched_sid>100001</if_matched_sid>
  <same_field>client_ip</same_field>
  <description>DoS detected: 50+ requests in 10s
    from same IP on ops-portal</description>
  <group>dos,attack,</group>
  <options>alert_by_email</options>
</rule>
\end{lstlisting}

The 50-request threshold was chosen to exceed realistic legitimate browsing
behaviour while remaining sensitive enough to detect a flood within its first
10-second window. The \texttt{same\_field} constraint ensures that rule 100002
counts only events from a single source IP, so traffic from multiple distinct
users does not accumulate into a false positive.

\subsection{Response}

Rule 100002 triggers the active response chain described in
Section~\ref{sec:monitoring}. The source IP is blocked at the
\texttt{DOCKER-USER} iptables chain within approximately 5 seconds of the
threshold being crossed. The block is automatically lifted after 300 seconds,
after which the attacker's requests are accepted again unless a new flood
triggers another block.

\subsection{Evaluation}

The detection mechanism functions correctly end-to-end. In testing, rule 100002
fires within the first 10-second window of a sustained flood, and the IP block
is applied before the second window begins. The Grafana \emph{DoS Alerts} stat
panel (rule 100002 counter) increments within seconds of the attack starting.

The principal limitation is that the threshold is calibrated for the simulated
load: a slow flood sustaining just below 50 requests per 10 seconds would not
trigger rule 100002. Additionally, single-IP blocking provides no defence
against a distributed attack where each source IP individually remains below
the threshold. That scenario is addressed separately in Section~\ref{sec:botnet}.
